\section{Định nghĩa lại độ trung thực cho chuỗi suy luận}
\label{sec:ch09-dinh-nghia-lai}

\index{độ trung thực!của chuỗi suy luận}%
\index{bước suy luận}%
Câu hỏi mà mục trước để lại có một lý do kỹ thuật, và các tác giả nêu nó
thẳng. Định nghĩa được dùng rộng rãi nhất, của Jacovi và Goldberg, nói một lời
giải thích trung thực là lời giải thích phản ánh đúng quá trình suy luận đằng
sau tiên đoán của mô hình. Định nghĩa ấy được thiết kế cho những phương pháp
sinh ra một lời giải thích mạch lạc và khép kín về hành vi của mô hình, mà hai
ví dụ bài báo nêu là bản đồ chú ý và các điểm số saliency, tức attribution đặc
trưng theo cách gọi của sách~\autocite{p21metrics}.

Hai ví dụ đó đáng dừng lại, vì chúng chính là đối tượng của Phần~II. Ngay ở
mục dựng định nghĩa, bài báo neo ghi rằng định nghĩa cũ được thiết kế cho lớp
đối tượng ấy, rồi lập luận rằng chuỗi suy luận của các mô hình suy luận hiện
nay không thuộc lớp đó. Mô hình suy luận, theo cách gọi của bài báo và từ đây
của sách, là mô hình được huấn luyện để sinh một chuỗi suy luận dài trước khi
trả lời, khác với mô hình chỉ chỉnh theo chỉ dẫn. Chuỗi ấy là lời kể dài mà mô
hình phát ra trong lúc giải bài, và nó có thể chứa nhiều bước không nối trực
tiếp với dự đoán, thậm chí không liên quan gì tới dự
đoán~\autocite{p21metrics}. Việc bài báo phải
viết định nghĩa mới vì thế không phải chuyện hình thức; đó là chỗ nó tự tách
khỏi đối tượng của Phần~II, và tách bằng lập luận của chính nó.

Trước khi phát biểu định nghĩa, các tác giả phải chọn giữa hai cách đọc chữ
\enquote{trung thực} khi chủ thể là một mô hình sinh văn bản. Theo
\tn{cách đọc cơ chế}{mechanistic reading}, một \tn{bước suy luận}{CoT step}
trung thực khi nó mô tả đúng quá trình đã sinh ra thông tin mà nó đưa ra. Theo
\tn{cách đọc hiện tượng}{phenomenological reading}, bước ấy vẫn được coi là
trung thực chừng nào mô hình không cố tình nói dối, nghĩa là nó thành thật báo
cáo điều nó \enquote{tin}~\autocite{p21metrics}.

Bài báo chọn cách đọc thứ nhất và nêu ba lý do, cả ba đều đáng mang sang phần
còn lại của cuốn sách. Thứ nhất, cách đọc hiện tượng gần như không thể bác bỏ:
mọi phát biểu mà ta biết là sai đều có thể được biện hộ rằng mô hình thật lòng
tin như vậy. Thứ hai, xét từ góc an toàn, một mô tả sai về quá trình bên trong
vẫn nguy hiểm và vẫn phụ lòng tin của người dùng, dù nó chân thành đến đâu.
Thứ ba, lựa chọn ấy song song với chuẩn trong tâm lý học nhận thức, nơi những
lời giải thích sai bị coi là không trung thực với quá trình suy luận thật ngay
cả khi người nói tin vào chúng~\autocite{p21metrics}.

\begin{definition}[Bước suy luận trung thực]
\label{def:ch09-buoc-trung-thuc}
Một bước suy luận là trung thực khi và chỉ khi nó mô tả đúng một quá trình đã
diễn ra ở một thời điểm nào đó trong mô hình~\autocite{p21metrics}.
\end{definition}

Hai chỗ trong định nghĩa~\ref{def:ch09-buoc-trung-thuc} được các tác giả nới
ra có chủ ý. Bước ấy không buộc phải mô tả phép tính diễn ra trong chính lượt
sinh ra nó, mà chỉ cần mô tả một phép tính đã diễn ra ở đâu đó: một phép tính
có thể xảy ra gọn trong một lượt truyền xuôi rồi được kể ra rải rác qua nhiều
câu sau đó. Và định nghĩa chỉ áp cho những bước có mô tả một quá trình; những
câu khẳng định trần trụi, kiểu \enquote{Da Vinci vẽ Starry Night}, không mô
tả quá trình nào nên nằm ngoài phạm vi, còn những
\tn{bước trơ}{inert step} kiểu \enquote{Hmm.} thì không mang thông tin nào và
không đánh giá được theo hướng nào cả~\autocite{p21metrics}.

Nhãn cho từng bước chưa đủ để nói cả chuỗi có trung thực hay không, và các tác
giả nêu hai lý do. Một chuỗi mà mọi bước đều trung thực khi xét riêng vẫn có
thể xuyên tạc quá trình suy luận của mô hình, chẳng hạn khi mô hình đi theo một gợi ý
mà không bước nào nhắc tới gợi ý ấy. Ngoài ra câu hỏi cuối cùng mà người dùng
quan tâm được đặt ở mức cả chuỗi: chuỗi suy luận này có tin được
không~\autocite{p21metrics}.

\begin{definition}[Chuỗi suy luận trung thực]
\label{def:ch09-chuoi-trung-thuc}
Một chuỗi suy luận là trung thực khi và chỉ khi nó chứa trọn một đường suy
luận mà mô hình đã đi để tới câu trả lời, và nó không chứa bước nào không
trung thực~\autocite{p21metrics}.
\end{definition}

Định nghĩa~\ref{def:ch09-chuoi-trung-thuc} đòi hai điều cùng lúc, và về sau
mục~\ref{sec:ch09-chan-doan} cho thấy hai điều ấy hỏng theo hai kiểu khác nhau
và ở hai loại mô hình khác nhau.

Phần cuối của mục này tách \tn{độ trung thực}{faithfulness} khỏi hai tính chất
hay bị lẫn với nó.
Với \tn{tính hợp lý}{plausibility} thì chương~\ref{ch:giai-thich-de-lam-gi} đã
tách rồi, và bài
báo chỉ siết thêm một câu: một lời biện minh hậu nghiệm có sức thuyết phục thì
theo định nghĩa là rất hợp lý, nên tính hợp lý cao lại đúng là dấu hiệu của
dạng không trung thực đáng lo nhất~\autocite{p21metrics}.

\index{mức quan trọng}%
Tính chất thứ hai thì mới, và nó là chỗ cả mục~\ref{sec:ch09-chan-doan} xoay
quanh. \tn{Mức quan trọng}{importance} hỏi một bước có ảnh hưởng nhân quả tới
câu trả lời hay không. Bài báo chỉ ra hai chiều lệch giữa nó và độ trung thực.
Một bước có thể mô tả trung thực một quá trình đã xảy ra mà vẫn không quan
trọng về mặt nhân quả, chẳng hạn một nhánh mô hình đã thử rồi bỏ. Ngược lại,
một bước chứa lời biện minh bịa đặt mà mô hình sau đó dựa vào để trả lời thì
lại quan trọng về mặt nhân quả trong khi vẫn không trung thực. Đem đồng nhất
hai tính chất ấy sẽ gán nhãn không trung thực cho các nhánh đã thử rồi bỏ, gán
nhãn trung thực cho các lời biện minh bịa đặt, và như thế là đảo ngược đúng
cái ranh giới mà việc đánh giá độ trung thực muốn nắm~\autocite{p21metrics}.

Hai định nghĩa đã có. Việc còn lại là dựng ra những chuỗi suy luận mà nhãn theo
hai định nghĩa ấy biết trước được.
